\section{Conclusion}

This project asks a simple question with a practical goal: how much speech is actually needed for an ALS speech assessment? Using VOC-ALS, I will compare every possible combination of the eight available smartphone speech tasks while keeping the evaluation participant-disjoint. The goal is not to build an ALS diagnostic tool, but to see whether remote speech monitoring can use a shorter protocol without losing much useful information.

The biggest limitation is the small dataset, especially for participants with more severe speech impairment. Task count is also only a rough estimate of burden and does not directly measure fatigue. Because of these limits, the final report will focus more on uncertainty and patterns that stay consistent across folds than on small score differences. If a shorter task set performs reliably close to the full protocol, it could be a useful starting point for lower-burden smartphone or assistive-technology assessments.

